%%%%%%%%%%%%%%%%%%%%%%%%%%%%%%%%%%%%%%%%%
% JPL Scientist II Research Statement
% Derived from main.tex in this directory.
%%%%%%%%%%%%%%%%%%%%%%%%%%%%%%%%%%%%%%%%%

\documentclass[11pt,letterpaper,sans]{moderncv}

\moderncvstyle{casual}
\moderncvcolor{grey}

\usepackage[scale=0.83]{geometry}
\usepackage[hidelinks]{hyperref}
\usepackage{enumitem}
\setlist[itemize]{leftmargin=1.2em,itemsep=0.2em,topsep=0.2em}
\setlength{\hintscolumnwidth}{2.0cm}
\setlength{\footskip}{41pt}

\firstname{Nhat-Minh}
\familyname{Nguyen}
\title{Research Statement: Scientist II, Roman and Euclid Cosmology}
\address{Kavli Institute for the Physics and Mathematics of the Universe}{University of Tokyo}
\email{nhat.minh.nguyen@ipmu.jp}
\homepage{https://minhmpa.github.io/}

\begin{document}
\makecvtitle
\vspace{-1.2em}

\section{Research vision}
My research asks the question, what representation of galaxy-survey data is best for the physics we want to test? For Roman and Euclid, the answer will not be one method or one summary statistic. Their weak-lensing and galaxy-clustering data will measure both geometry and growth with high statistical power. The hard part is to interpret those data without letting nonlinear structure formation, galaxy bias, intrinsic alignments, selection effects, or likelihood approximations masquerade as new physics.

I develop field-level and forward-modeling methods for galaxy surveys. My niche is likelihood validation and information accounting---choosing the data representation that is optimal for a given physical question, then testing where the chosen likelihood loses information or becomes biased. At the Jet Propulsion Laboratory (JPL), I would contribute to Roman and Euclid weak-lensing and clustering analyses through nonlinear modeling, likelihood validation, and systematic-controlled tests of growth, gravity, and dark energy.

\section{Preparation: methods and science background}
My background combines the science question with the method needed to make the answer trustworthy.

On the science side, my first-author Physical Review Letters (PRL) paper reported evidence for a late-time suppression in the growth of large-scale structure. That work used growth-rate information from the cosmic microwave background (CMB), weak lensing, galaxy clustering, and cosmic velocities to test the consistency of the concordance cosmological model, and it led naturally into modified-gravity interpretations with collaborators. It shaped how I think about Roman and Euclid. Their measurements can connect expansion history to growth history to test the standard models of gravity and cosmology, but that connection is only as strong as the modeling and systematics control behind it.

On the methodology side, field-level inference is my main technical contribution. I have developed and validated Bayesian forward-modeling approaches that infer cosmology from the observed galaxy field itself. In a first-author PRL with collaborators, I used an effective-field-theory (EFT) differentiable forward model, LEFTfield, to quantify how much information nonlinear galaxy clustering contains beyond low-order summary statistics. This work was recognized by the 2024 Buchalter Cosmology Prize. For Roman and Euclid, the value is not that every analysis should become field-level. The value is that field-level methods give a controlled way to test information loss, validate likelihoods, stress-test scale cuts, and decide where summary statistics are sufficient.

I have also worked on survey-facing problems that connect directly to the JPL role. In the Prime Focus Spectrograph (PFS), I am responsible for the full-shape galaxy-clustering analysis pipeline. My single-author 2026 arXiv preprint, submitted to JCAP, ``Multi-tracers, multi-surveys: a joint Fisher analysis of DESI+PFS,'' grows out of that role and studies how overlapping surveys and multiple galaxy populations can calibrate nuisance parameters in full-shape analyses for PFS and the Dark Energy Spectroscopic Instrument (DESI). My earlier work on field-level inference of intrinsic alignment from Sloan Digital Sky Survey Baryon Oscillation Spectroscopic Survey galaxy samples connects galaxy shapes, tidal fields, and weak-lensing systematics. This is the contribution I would bring to Roman and Euclid: connect calibrated shape catalogs, photometric-redshift information, clustering data, and mission mocks to model likelihood validation and dark-sector theory interpretation.

\section{Proposed JPL contribution}
I would focus my JPL contribution on validation and interpretation for joint weak-lensing and clustering cosmology.

First, I would contribute to joint modeling of Roman and Euclid observables. Weak lensing constrains the projected matter field; galaxy clustering traces biased tracers in three dimensions. A useful analysis must model galaxy bias, matter clustering, intrinsic alignment, photometric-redshift leakage into cross-correlations, source--lens coupling, selection, and covariance together rather than treating them as unrelated nuisance terms. I would work on diagnostics that sit downstream of existing shear, photo-z, and clustering products: which combinations of lensing and clustering measurements give the most robust growth and gravity tests, and where cross-covariance with clustering controls the error budget.

Second, I would build simulation-based likelihood checks. This is where my field-level background is most useful. Using mission mocks and OpenUniverse-style simulations where appropriate, I would compare standard two-point analyses, higher-order summaries, hybrid compression, and field-level or forward-modeled likelihoods under apple-to-apple matched assumptions. The validation target would be practical: scale-cut sweeps, covariance and cross-covariance checks, prior-sensitivity tests, null tests, and documented comparisons showing where standard likelihoods are reliable, where they compress (and lose) information, and where model error or systematics dominate.

Third, I would connect the validated measurements to dark-sector interpretation. Roman and Euclid will sharpen tests of whether geometry and growth agree. My growth-suppression and modified-gravity work gives me a foundation for interpreting those tests without over-reading them. I would aim to turn modeling results into constraints on growth, gravitational slip, dark-energy evolution, and LambdaCDM consistency, with uncertainty budgets that make the limiting assumptions visible.

\section{Three-year integration and independent niche}
As a Scientist II, I would enter JPL as a focused contributor with an independent technical niche: field-level and forward-modeling validation for Roman/Euclid dark-sector inference. In the first year, I would integrate with the relevant Roman and Euclid working groups, learn the current likelihood and simulation infrastructure, and contribute to an existing weak-lensing or clustering modeling task. In parallel, I would identify one validation problem where my background can add value quickly, such as a Roman--Euclid overlap consistency test, an intrinsic-alignment or galaxy-bias stress test, or a scale-cut comparison across summary-statistic and forward-modeled analyses.

In years one to two, I would aim to lead, within the relevant working group, a focused analysis paper or validation product. The product should be useful even if the final mission analysis remains summary-statistic based: a Roman--Euclid overlap mock challenge, a likelihood-comparison note, or model-error diagnostics usable by the working group. In years two to three, I would expand that work into a focused independent niche connecting weak lensing, clustering, and field-level information accounting.

I would bring a collaborative working style to this role. I have worked in PFS and DESI contexts, contributed to community validation efforts, chaired the organization of the Beyond-2-Point Statistics Meet Survey Systematics workshop, and convened the large-scale-structure session at COSMO'24. I have also advised and co-mentored students and early-career researchers whose projects led to papers or accepted workshop publications: Andrja Kosti\'c's EFT-likelihood consistency work became a JCAP paper, and Baptiste Barthe-Gold's density-field reconstruction project was accepted to the NeurIPS 2025 Machine Learning for Physical Science Workshop.

That is also how I tend to work. I learn unfamiliar territory by building the first usable version of an analysis, then making it clear enough that someone else can test or improve it. At JPL, I would use that habit in service of mission teams: precise enough for experts, documented enough for collaborators, and cautious enough that any claimed tension with LambdaCDM or general relativity is supported by the modeling.

Roman and Euclid need nonlinear modeling, systematics control, and inference methods that can connect weak lensing and galaxy clustering to fundamental physics. I would use my background to reinforce the robustness and interpretability of thoese connections.

\end{document}
